\section{Methodology}

\subsection{Datasets and Evaluation}
We use the Arabic Sign Language dataset (ArASL)~\cite{latif2020arsl}, a public
dataset of over 54{,}000 grayscale images covering the 32 Arabic alphabet signs,
collected from more than 40 contributors under varied hand positions.

A practical issue with this dataset is that it contains many near-duplicate frames
of the same person holding the same letter. If the images are split randomly,
almost-identical frames end up in both the training and test sets, so the model can
succeed on test images simply because it saw their near-twins during training. This
is data leakage, and it makes test accuracy look much higher than the model's real
ability. Our first split was random, and the resulting accuracy turned out to be
inflated for exactly this reason.

To fix this, we use a leakage-aware split. The dataset has no signer labels, so we
approximate one: we compare images with a perceptual hash (dHash), group
near-duplicate frames into clusters, and send each cluster as a whole to only one of
the training, validation, or test sets (70/15/15). After clustering, the RGB images
formed 7{,}575 clusters out of 7{,}856 images; assigning whole clusters gives
5{,}509 training, 1{,}182 validation, and 1{,}165 test images, with no cluster
shared across any two splits. This guarantees that no frame and its near-copies can
appear on both sides of the split.

To check generalization beyond one source, we also use a second, independent
dataset: the RGB Arabic Alphabets Sign Language Dataset, which is in color and was
collected separately. After matching the class names between the two datasets (they
use different spellings, and the RGB set has no \texttt{yaa}, which we drop), we
obtain 31 shared classes. We evaluate in two ways: a clean held-out test set from
the same distribution, and a cross-dataset test on the independent set.

\subsection{Preprocessing and Data Augmentation}
All images are resized to $224 \times 224$ to match the MobileNetV2 input, turned
into tensors, and normalized with the standard ImageNet statistics. We also focus
the model on the hand region using MediaPipe-based detection and cropping, which
removes much of the background variation that hurts performance on new data. The
same preprocessing is used during training and during live inference, so the model
sees images the same way in both cases.

During training we apply light augmentation to imitate real webcam conditions: small
random rotations (up to about $\pm10^\circ$), small affine shifts and scaling, and
mild brightness and contrast changes. We deliberately do not use horizontal
flipping, because mirroring a hand can turn one sign into a different one.
Validation and test images are not augmented.

\subsection{Model Architecture}
We use MobileNetV2~\cite{sandler2018mobilenetv2} as the backbone. It is a good fit
here because it is accurate while staying light enough to run in real time on a
normal laptop. Instead of training from scratch, we start from ImageNet pretrained
weights~\cite{deng2009imagenet} and adapt them to our task. The final classification
layer is replaced to output the number of sign classes (32 for the ArASL model, and
31 for the domain-adapted RGB model after dropping \texttt{yaa}).

Training has two phases. First we train only the new classifier layer while the
backbone is frozen, so the network can adjust to the new labels without damaging the
pretrained features. Then we unfreeze the whole network and fine-tune it at a smaller
learning rate. All training is done in PyTorch~\cite{paszke2019pytorch} on a Kaggle
NVIDIA Tesla T4 GPU.

\subsection{Domain Adaptation to Color Data}
After the leakage analysis, we adapt the model to the color (RGB) distribution,
which is closer to what a real webcam produces. We initialize from the ImageNet and
ArASL weights and continue training on the RGB dataset using the leakage-aware split
described above. Because the test set here is genuinely unseen, the resulting
accuracy is a fair estimate of how the model performs on new data, rather than an
inflated benchmark number.

\subsection{Hand-Aware Background Blurring}
To make the live interface easier to read, we blur the background around the
detected hand. The steps are:
\begin{enumerate}
  \item Detect 21 hand landmarks with MediaPipe Hands~\cite{lugaresi2019mediapipe}.
  \item Build a soft elliptical mask $M$ centered on the hand, with edges smoothed by
        a Gaussian blur of the mask.
  \item Blur the whole frame with a Gaussian kernel ($51 \times 51$) to get a blurred
        copy $I_{\text{blur}}$.
  \item Blend the sharp and blurred frames using the mask:
        \begin{equation}
          I_{\text{out}} = I_{\text{sharp}} \cdot M + I_{\text{blur}} \cdot (1 - M)
        \end{equation}
\end{enumerate}
The result keeps the hand sharp while softening the background, and because the mask
edges are smoothed, the transition looks natural. The same idea could later be used
to mask the background before classification, which we leave for future work.

\subsection{Training Strategy}
We train with the AdamW optimizer and cross-entropy loss with label smoothing
($\epsilon = 0.10$), which discourages overconfident predictions. The classifier head
is trained for 5 epochs with the backbone frozen, then the full network is fine-tuned
for up to 25 epochs at a lower learning rate. We use ReduceLROnPlateau to lower the
learning rate when validation stops improving, mixed-precision training (AMP) for
speed, and early stopping (patience 6). The best model on the validation set is kept.
The full configuration is shown in Table~\ref{tab:training-config}.

\textit{Note on the loss values:} with label smoothing, the cross-entropy loss has a
non-zero floor (around $0.37$ for this number of classes), so a loss near that value
at high accuracy is expected and does not mean the model is underfitting.